% =============================================================
%  CV - Ha Van Doan
%  BIEN DICH BANG XeLaTeX (khong dung pdfLaTeX - se loi tieng Viet)
%
%  Overleaf : Menu > Compiler > XeLaTeX > Recompile
%  Local    : xelatex cv.tex
%
%  Chi tiet xem README.md trong cung thu muc nay.
% =============================================================
\documentclass[11pt,a4paper]{article}

\usepackage{iftex}
\ifXeTeX\else\ifLuaTeX\else
  \GenericError{}{Tai lieu nay phai bien dich bang XeLaTeX hoac LuaLaTeX}{}{}
\fi\fi

% ---------- fonts ----------
% TeX Gyre di kem moi ban TeX Live / MiKTeX / Overleaf va co day du dau tieng Viet.
\usepackage{fontspec}
\defaultfontfeatures{Ligatures=TeX}
\setmainfont{texgyreheros}[
  Extension      = .otf,
  UprightFont    = *-regular,
  BoldFont       = *-bold,
  ItalicFont     = *-italic,
  BoldItalicFont = *-bolditalic,
]
% Muon CV kieu chan (serif) thi doi 'texgyreheros' thanh 'texgyretermes'.

% ---------- layout ----------
\usepackage[a4paper,top=1.4cm,bottom=1.4cm,left=1.6cm,right=1.6cm]{geometry}
\usepackage{xcolor}
\usepackage{titlesec}
\usepackage{enumitem}
\usepackage{array}
\usepackage{tabularx}
\usepackage{hyperref}

\definecolor{accent}{HTML}{0A7A63}
\definecolor{ink}{HTML}{101826}
\definecolor{muted}{HTML}{556579}
\definecolor{rule}{HTML}{D6DAE0}   % xam trung tinh, khong nga xanh

\color{ink}
\hypersetup{
  colorlinks = true,
  urlcolor   = accent,
  pdfborder  = {0 0 0},
  pdftitle   = {CV - Ha Van Doan},
  pdfauthor  = {Ha Van Doan},
}

\pagestyle{empty}
\setlength{\parindent}{0pt}
\linespread{1.06}

% ---------- section heading ----------
% Duong ke duoi tieu de: mot doan ngan mau accent noi tiep bang mot net toc xam.
\newcommand{\sectionrule}{%
  \color{accent}\rule[0.5pt]{1.4cm}{1.6pt}%
  \color{rule}\rule[1pt]{\dimexpr\linewidth-1.4cm\relax}{0.6pt}%
}

\titleformat{\section}
  {\large\bfseries\color{accent}}
  {}{0pt}
  {\MakeUppercase}
  [\vspace{-7pt}{\sectionrule}]
\titlespacing*{\section}{0pt}{13pt}{8pt}

% ---------- icons (tu dong bo qua neu khong co fontawesome5) ----------
\IfFileExists{fontawesome5.sty}{%
  \usepackage{fontawesome5}
  \newcommand{\icoMail}{\faEnvelope[regular]~}
  \newcommand{\icoGit}{\faGithub~}
  \newcommand{\icoPhone}{\faPhone*~}
  \newcommand{\icoPin}{\faMapMarkerAlt~}
}{%
  \newcommand{\icoMail}{}
  \newcommand{\icoGit}{}
  \newcommand{\icoPhone}{}
  \newcommand{\icoPin}{}
}

% =============================================================
%  MACRO
% =============================================================

% \role{chuc danh}{to chuc}{thoi gian}
\newcommand{\role}[3]{%
  \vspace{5pt}\par
  \begin{tabularx}{\linewidth}{@{}X r@{}}
    \textbf{#1}\ \textcolor{muted}{\textbar}\ #2 & \textcolor{muted}{\small #3}
  \end{tabularx}\par
  \vspace{2pt}%
}

% \project{ten}{nhan}{link github}{cong nghe}
\newcommand{\project}[4]{%
  \vspace{6pt}\par
  \begin{tabularx}{\linewidth}{@{}X r@{}}
    \textbf{#1}\ \textcolor{muted}{\textbar}\ \textcolor{muted}{\small #2}
    & \href{#3}{\icoGit\ttfamily\footnotesize source}
  \end{tabularx}\par
  {\footnotesize\color{accent}#4\par}%
  \vspace{2pt}%
}

% bullet list gon
\newlist{pts}{itemize}{1}
\setlist[pts]{
  leftmargin = 12pt,
  labelsep   = 5pt,
  topsep     = 2pt,
  itemsep    = 1pt,
  parsep     = 0pt,
  label      = \textcolor{accent}{\small\textbullet},
}

% =============================================================
%  NOI DUNG
% =============================================================
\begin{document}

% ---------- HEADER ----------
{\fontsize{26}{30}\selectfont\bfseries HÀ VĂN ĐOÀN}\\[4pt]
{\large\color{accent}Flutter / Mobile Developer\ \textcolor{muted}{\textbar}\ AI \& IoT}\\[7pt]
{\small
  \icoMail\href{mailto:bestbubuom@gmail.com}{bestbubuom@gmail.com}
  \quad\textcolor{rule}{\textbar}\quad
  \icoGit\href{https://github.com/Hdoanf}{github.com/Hdoanf}
  % --- BO SUNG: bo dau % o dau dong va dien thong tin cua ban ---
  % \quad\textcolor{rule}{\textbar}\quad \icoPhone 0xxx xxx xxx
  % \quad\textcolor{rule}{\textbar}\quad \icoPin Hà Nội, Việt Nam
}

\vspace{5pt}
{\color{rule}\rule{\linewidth}{0.6pt}}

% ---------- TOM TAT ----------
\section{Tóm tắt}

Sinh viên CNTT, tập trung vào \textbf{Flutter} và các bài toán mà ứng dụng phải kết nối với
thế giới bên ngoài: cảm biến, mô hình AI, hoặc kho dữ liệu. Đã hoàn thành 3 ứng dụng di động
đầy đủ chức năng, trong đó đồ án tốt nghiệp tự làm trọn vòng từ firmware \textbf{ESP32-C3}
$\to$ \textbf{BLE} $\to$ app Flutter $\to$ \textbf{Firebase} $\to$ \textbf{Gemini AI}.
Thành thạo dùng \textbf{AI agent} (Claude Code, Gemini CLI, MCP) như một phần quy trình
code hằng ngày.

% ---------- KY NANG ----------
\section{Kỹ năng}

\begin{tabularx}{\linewidth}{@{}>{\raggedright\arraybackslash}p{3.2cm} X@{}}
  \textbf{Mobile}          & Flutter, Dart, Riverpod, Provider, go\_router, fl\_chart, flutter\_tts, Responsive UI \\[3pt]
  \textbf{AI agent}        & Claude Code, Gemini CLI, MCP, prompt engineering, context engineering, AI code review \\[3pt]
  \textbf{AI / ML}         & Google Gemini, LangChain, FAISS, RAG, ML Kit OCR, pypdf \\[3pt]
  \textbf{IoT / Embedded}  & ESP32-C3, BLE, cảm biến MAX30102, MJPEG stream, C / C++ \\[3pt]
  \textbf{Backend}         & Firebase Auth, Cloud Firestore, FastAPI, Python, REST API \\[3pt]
  \textbf{Web (đồ án)}     & PHP, MySQL, HTML/CSS, JavaScript, PHPMailer, PHPWord \\[3pt]
  \textbf{Công cụ}         & Git, Neovim / LazyVim, Linux, Nix, Composer, CMake \\
\end{tabularx}

% ---------- KINH NGHIEM ----------
\section{Kinh nghiệm}

\role{Thực tập sinh Flutter}{Công ty Cổ phần HC-Hiteck}{12/2025 -- 04/2026}
\begin{pts}
  \item Xây dựng ứng dụng quản lý \textbf{Smart Home IoT} đa nền tảng --- một codebase chạy
        được cả mobile lẫn desktop với layout responsive riêng cho từng kích thước màn hình.
  \item Tích hợp xem camera an ninh \textbf{MJPEG} realtime, hệ thống cảnh báo cháy và
        bộ hẹn giờ tự động bật/tắt thiết bị.
  \item Trực quan hoá dữ liệu sử dụng thiết bị bằng biểu đồ tương tác (FL Chart);
        quản lý state bằng \textbf{Riverpod}, điều hướng bằng \textbf{go\_router}.
  \item Tổ chức mã nguồn theo kiến trúc \textit{feature-first} (auth, device, fire\_alert,
        scheduler, stats\ldots) để dễ mở rộng và bàn giao.
\end{pts}

% ---------- DU AN ----------
\section{Dự án nổi bật}

\project{HealthPulse VN}{Đồ án tốt nghiệp}{https://github.com/Hdoanf/doantotnghiep}%
{Flutter · Firebase · Google Gemini · ML Kit OCR · ESP32-C3 · BLE}
\begin{pts}
  \item Ứng dụng quản lý sức khoẻ cá nhân kết hợp \textbf{AI} và \textbf{IoT}.
  \item Tự viết firmware cho \textbf{ESP32-C3} kèm cảm biến \textbf{MAX30102}, truyền nhịp tim
        và SpO\textsubscript{2} về app qua \textbf{BLE} theo thời gian thực.
  \item Chat tư vấn sức khoẻ dựa trên chỉ số người dùng bằng \textbf{Google Gemini};
        quét phiếu xét nghiệm giấy bằng \textbf{ML Kit OCR}.
  \item Biểu đồ theo dõi chỉ số sinh tồn và hồ sơ sức khoẻ lưu trên Cloud Firestore.
\end{pts}

\project{App Học Tiếng Trung}{Sản phẩm cá nhân}{https://github.com/Hdoanf/apphoctiengtrung}%
{Flutter · Firebase Auth · Firestore · flutter\_tts · Custom Painter}
\begin{pts}
  \item Canvas viết tay chữ Hán kèm \textbf{bộ đánh giá nét tự viết} (có unit test riêng)
        dựa trên dữ liệu thứ tự nét của từng chữ.
  \item Luyện nghe--nói bằng TTS, bài tập ghép câu, theo dõi tiến độ người học trên Firestore.
  \item Tự thiết kế design system riêng: bảng màu, typography và widget dùng chung.
\end{pts}

\project{Chatbot AI Amy --- Kho tài liệu sinh viên}{RAG}{https://github.com/Hdoanf/demo}%
{Python · FastAPI · LangChain · FAISS · pypdf · boto3}
\begin{pts}
  \item Chatbot trả lời câu hỏi dựa trên chính tài liệu người dùng upload, dùng pipeline
        \textbf{RAG} với LangChain và vector store \textbf{FAISS}.
  \item Trích xuất nội dung PDF bằng pypdf, lưu file lên S3 qua boto3, backend FastAPI.
\end{pts}

\project{chimmusicplayer}{Open source}{https://github.com/Hdoanf/chimmusicplayer}%
{C · TUI · Linux · Makefile · Nix}
\begin{pts}
  \item Trình phát nhạc chạy trong terminal viết bằng \textbf{C}, tuỳ biến từ \texttt{kew}.
  \item Bổ sung tìm kiếm YouTube ngay trong ứng dụng, tự tải MP3 về máy và nhúng ảnh bìa
        cùng metadata bài hát vào file để nghe offline.
\end{pts}

\project{Hệ thống quản lý tài sản \& kho}{Đồ án môn học}{https://github.com/Hdoanf/hdoanf.github}%
{PHP · MySQL · PHPWord · php-barcode-generator}
\begin{pts}
  \item Quản lý tài sản, kho và người dùng có phân quyền; sinh mã vạch cho từng tài sản
        và xuất báo cáo ra file Word.
\end{pts}

% ---------- HOC VAN ----------
\section{Học vấn}

\role{Công nghệ thông tin}{Đại học Hòa Bình}{2022 -- 2026}
\begin{pts}
  \item GPA: 3.3 / 4.0
  \item Đồ án tốt nghiệp: \textbf{HealthPulse VN} --- ứng dụng sức khoẻ tích hợp AI và IoT.
\end{pts}

\end{document}
